\section{Overview of the construction verifying (CP0)--(CP8)}\label{sec:design}

\paragraph{What this section proves and how to read it.}
Part~V exhibits laws satisfying (CP0)--(CP8) with baseline
$\Theta<\PushTheta<1$. This section specifies the construction once:
parameters, queries, prime ranges, and which graph law is used in
each range. Later sections verify the hypotheses range by range;
each ends by naming the rows of \cref{tab:hypotheses} it proves.
Inputs: Theorem~C and Part~IV. Output: the standing data
\cref{def:record:parameters} for the baseline. The improved
($\vartheta>1$) parameters differ and are specified in
\crefrange{sec:density:moments-prefix}{sec:four-routes-closing};
do not use baseline numbers there. ``Admissible'' in Part~V has
three qualified senses, distinguished in
\cref{subsec:design-admissibility}: admissible for the energy
estimates (\cref{def:admissible}), admissible tuples (local degree
conditions), and admissible price bounds $c_{\mathcal C}$ (allowed
(CP5) scalars). Where a sum is computed at a smaller evaluation
exponent and converted to $a$ (repricing),
\cref{lem:record:repricing,cor:first-stage-repricing} justify the step
and require $a\ge a_0$.

The prime ranges are: prefix $p\equiv1\pmod4$, $p<2^{512}$
(\cref{sec:frontier:iidprefix}); middle $p\equiv3\pmod4$ in
$[Q_3,P_3)$ (\crefrange{sec:middle-reference}{sec:matching-mean});
upper $p\equiv3\pmod4$ in $[P_3,2^{512})$
(\cref{sec:upper-reference,sec:upper-comparison}); large primes
$p\ge2^{512}$ and the external index (\cref{sec:tail-index}); closing
(\cref{sec:closing}). Baseline totals are
$L_0,S,A,D,\mathcal E,T,R_I,J,W,m$ as in
\cref{subsec:intro:inequality,tab:intro:values}.

\subsection{The baseline numerical parameters}\label{subsec:design-parameters}

The construction is now fixed completely: exponents, number of
intervals, prime ranges, and which law is used in each range.
These are the baseline values ($\theta=\vartheta=1$); the improved
closing changes only the values named in
\cref{sec:four-routes-closing}.

\begin{definition}[Parameters]\label{def:record:parameters}
The construction has the following data.
\begin{enumerate}[label=(\alph*)]
\item \emph{Exponents.} The price exponent $a=\PushA$, the density exponent
 $\theta=1$, and the exponent of the baseline theorem
 $\eps=a/(72-\theta)=a/71=\PushExponent$.
\item \emph{Slots, queries and intervals.} The slot array
 $V=(\Zsix)^2\times\{0,1\}$ with its family $\mathfrak F$ of cuts and its
 reserved cut $\iota$ (\cref{def:slot-array}); the six queries $\mathcal E_6$
 (\cref{subsec:front6:construction}); and $\Kint=\PushK$ intervals.
\item \emph{Prime ranges.}
 \begin{itemize}
 \item \emph{Prefix:} the primes $p\equiv1\pmod4$ below $\Pfour=2^{512}$. Each
  is an atom $\{p\}$ with label group $\Z/p^{e_p}\Z$, $e_p=2\mcl$, $\mcl=8192$ digit pairs, carrying the
  square-clique law of \cref{def:clique-law}.
 \item \emph{Small modulus:} $\Mzero=2^{2m_{\rm ext}+1}\prod_{3\le p<Q_3,\
  p\equiv3\ (4)}p^{2m_{\rm ext}}$, $m_{\rm ext}=64$, $Q_3=\FurtherQ$. Its primes are not coordinates of the
  label group; they enter only through the index (\cref{def:hyb:coupled}).
 \item \emph{Middle range:} the primes $p\equiv3\pmod4$ in $[Q_3,\Pthree)$,
  $\Pthree=\FurtherPthree$, in $33$ pieces $[X_i,X_{i+1})$
  (\cref{tab:record:middle-pieces}), each partitioned into exact blocks
  (\cref{def:record:exact-blocks}); the seven tilted pieces start at $2^{30}$, $2^{31}$, $2^{32}$,
  $2^{34}$, $2^{36}$, $2^{40}$ and $2^{45}$.
 \item \emph{Upper range:} the primes $p\equiv3\pmod4$ in $[\Pthree,\Pfour)$, in
  $310$ bands, each partitioned into blocks (\cref{def:record:upper-bands}):
  $45$ bands in $[\Pthree,2^{48})$, namely $[167\cdot2^{40},168\cdot2^{40})$
  and the $44$ bands $[j2^{41},(j+1)2^{41})$, $84\le j<128$; $64$ equal bands
  in each octave $[2^b,2^{b+1})$, $48\le b\le51$; and the nine bands
  $[2^b,2^{b+1})$, $52\le b\le57$, $[2^{58},2^{60})$, $[2^{60},2^{64})$ and
  $[2^{64},\Pfour)$.
 \item \emph{Tail:} every prime $p\ge\Pfour$, an atom $\{p\}$ with $e_p=1$; the
  threshold in the mixing condition \ref{it:record:R4}
 (\cref{lem:record:tail-mixing}) is $\Pfive=2^{65536}$.
 \end{itemize}
 Every other prime ($2$, and the primes $p\equiv3\pmod4$ below $Q_3$) divides
 $\Mzero$.
\item \emph{Block budgets.} Every middle block has $\log q_{\Blk}\le L=570$, so
 that its selection factor is at most $e^{2aL}$, with
 $2aL=1236349836/1220703125<1.013$. The local laws of the middle range are
 evaluated for blocks of logarithmic length at most the local cap
 $L_{\rm loc}=585$, at which their constants are computed
 (\cref{lem:record:local-domain}). Every upper block has $\log q_{\Blk}$ at
 most the recorded budget $L_{\rm band}$ of its band, an integer between $550$
 and $886$ (\cref{def:record:upper-bands}), at which the local data of the band
 are evaluated; the completed-root caps of the bands from $2^{52}$ and the
 local data of the tail law (\cref{lem:tail-law}) are evaluated for the larger block length
 $L_{\rm u}=50000/27$.
\item \emph{Fold records.} In the comparison law the two-layer fold records
 $F_1,F_2$ of the middle blocks are replaced by the signed whole-block fold
 records with parameter $\tlaw=\FurtherT$ (\cref{def:record:mixture}).
\item \emph{Roots.} On a middle piece starting below $2^{30}$ the
 two-layer root is the uniform strict-square labelling law
 $\lambda_{\Blk}(G_{\rm L,2})$. The other seven pieces use the positive
 reweightings proved in \cref{sec:positive-roots}:
 \begin{center}
 \begin{tabular}{lll}
 piece start & multiplier & $\ulaw$\\ \hline
 $2^{30}$ & $\frac13\sum_j(1+\ulaw\Hs j)^2$ & $711/4096$\\
 $2^{31}$ & same & $1021/4096$\\
 $2^{32}$ & same & $1455/4096$\\
 $2^{34}$ & same & $833/1024$\\
 $2^{36}$ & same & $3851/4096$\\
 $2^{40}$ & same & $4049/4096$\\
 $2^{45}$ & $1/100+(1+\ulaw H_{39})^2$ & $67/40$
 \end{tabular}
 \end{center}
 Each multiplier is normalized at every prime; $H_{39}=H_{\mathcal D_{39}}$ is the collected signed trace sum with coefficients
 $\pm1/15$ specified in \cref{sec:positive-roots}.
 The one-layer root remains $\lambda_{\Blk}(G_{\rm L,1})$ (the graph
 $G_{\rm L,1}$ of \cref{subsec:push:layer12spectral}) on every piece.
\item \emph{First stage and bracket.} The one-layer folded copies of the fold records are fully
 polarized (polarized fraction one, i.e.\ the whole of $F_{\rm fold}$ is
 subtracted in forming $\rho$; \cref{def:construction}(c)); the folded
 matching records use the
 one-stage repair of \cref{lem:record:one-stage-folded-e}, whose record
 differences, of the two types $E_0$ and $E_1$, form the two tags of (CP2)(b)
 (\cref{def:record:frame}(e)); the dominator $\Nbr$ of (CP3)
 carries one folded copy of each fold record
 (\cref{def:record:frame}(d)).
\item \emph{Closing constants.} The common-mean coefficient $63/64=1-2^{-6}$
 (\cref{thm:comparison} with $k=6$), and the side weight
 $\tside=31/32$ for the auxiliary sides of the joint Gram block
 (\cref{def:side-function}); every numerical side bound of
 the matching envelopes is at least the corresponding auxiliary side
 bound (\cref{prop:signed-sides}).
\item \emph{Credits.} Variance credits of degrees one to four, with positive
 backgrounds of degrees one and two (\cref{sec:variance}); mean savings of all
 degrees up to four (\cref{prop:record:mean-credit}); on the pieces starting at
 or above $2^{31}$, pooling over those pieces adds background
 degrees three and four and mean degrees five and six
 (\cref{lem:pool-coupons}); every restoration, i.e.\ every recovery of
 product energy as in \cref{subsec:variance:restoration}, is
 formed at the exponent $a$ with the piece product caps of \cref{lem:piece-caps}.
\item \emph{Analytic stage} (\cref{def:record:analytic-stage}). $\beta=1/20$,
 the major-arc exponent $\eta=17/100$, $\gamma=255/1024$, $\tau=1/201$,
 $t_{\rm an}=1/60$, $\sigma=1/825$, and the moment losses $\delta=10^{-10}$ and
 $\ell_0=10^{-8}$.
\end{enumerate}
\end{definition}

The completion weights of the $33$ middle pieces and of the $310$ upper bands,
and the other finite data of the construction, are listed or referenced in
\cref{app:data}; the finite consequences of these choices that the frame uses
(the tiling of $[Q_3,\Pthree)$ by the pieces, the completion inequalities and
the protection inequality \eqref{eq:protection} in every piece, the
one-stage gates, the band starts and the
prefix gate) are collected in \cref{cert:record:frame-data}
(\cref{app:certificates}). The price exponent is below the error exponent of
the analytic stage, $\sigma-a=65195254/201416015625>0$
(\cref{cert:record:frame-data}(a)); the strict inequalities (S1)--(S7) of the
analytic stage hold (\cref{cert:record:analytic-stage}), and the gates of the
analytic transfer that depend on $a$ hold on the whole interval
$0<a\le9/10000$, for every fixed $\Kint\ge2$
(\path{research/advance_2026_10_24/global/analytic_domain.py}; the baseline
verifier asserts $0<a\le9/10000$). The larger parameter domain used for
the main theorem is verified in \cref{sec:four-routes-closing}.
All data are fixed integers or rationals,
independent of the scale; the analysis of Part~II applies with constants that
depend on them (\cref{rem:fixed-parameters}). The certificate computes the
prices of the middle range at the anchor $a_0=\PushAnchor<a$ and reprices them by
the factor $t=e^{2(a-a_0)L}$, about $1.0000000616512$; it computes the prices of the
upper bands at the anchor $a_{\rm u}=883/10^6<a$ and reprices them band by band
with the logarithmic-product cap of each band
(\cref{lem:record:repricing,cor:first-stage-repricing}). The numbers $a_0$ and
$a_{\rm u}$ are points of evaluation of the certificate, not parameters of the
construction.

\subsection{The six queries \texorpdfstring{$\mathcal E_6$}{E\_6} and the twelve pairs in projections}\label{subsec:front6:construction}

The comparison law accepts the six queries
$e^{(j)}=((0,j,0),(\rot0,\rot j,1))=((0,j,0),(2,\rot j,1))$, $j\in\Zsix$, of
\cref{def:slot-array}(e), the family $\mathcal E_6$. By \cref{lem:slot-array}(iv)
their twelve endpoints are distinct, every source lies in $V_+=\{\ell=0\}$ and
every reflected target $\iota(2,\rot j,1)=(2,\rot j,0)$ lies in $V_+$; so
$\mathcal E_6$ is a family of $k=6$ queries in the sense of \cref{def:cp0}(c).
Each query crosses the reserved cut from layer $0$ to layer $1$ and applies
$\rot$ to both chamber coordinates.

All six queries have the same first-coordinate pair $(0,\rot0)=(0,2)$; they
differ only in the second chamber coordinate, on which they act by the
permutation $\rot$. A law that ignores the second coordinate therefore sees
one query. The middle laws are such laws: they live on one chamber coordinate
and the layer, and the pullback reads every query of $\mathcal E_6$ through
the single pair $((0,0),(2,1))$ (\cref{lem:slot-array}(v)). In the middle
records the six queries thus impose one directed edge between the layers: the
two-layer root and the aligned fold $F_0$ accept it (on the pieces from
$2^{30}$ the root is the positive root of \cref{def:positive-root}, which has
the accepting support of the uniform root), the two-layer fold records
$F_1,F_2$ and matching records $E_0,E_1$ need one added cross edge each, and
the folded matching records one edge between their two private labels (the two labels not shared
between its halves; \cref{cert:record:middle-geometry}). The square-clique law of the prefix
ignores the layer as well, and sees the six queries as one coincidence event
of two labels (\cref{lem:frontier:iidlocal}(v)). This is why the six queries
impose on the middle range the same single edge condition as one query,
while the matching term of the closing
carries the sixth power $(R_I/2)^6$ (\cref{thm:comparison}).

\begin{remark}[The twelve-query superset]\label{subsec:front12:construction}
The family
\[
 \mathcal E_{12}=\bigl\{\bigl((u,j,0),(\rot u,\rot j,1)\bigr):u\in\{0,5\},\
 j\in\Zsix\bigr\}
\]
consists of twelve queries with twenty-four distinct endpoints, each running
from layer~$0$ to layer~$1$. Its first-coordinate pairs are $(0,2)$ and
$(5,\rot5)=(5,3)$, and $\mathcal E_6$ is its part with $u=0$. Some parts of the
construction accept all twelve queries: the one-way repairs of
the upper range, whose ledger was built for $\mathcal E_{12}$
(\cref{sec:frontier:upper-oneway}), the dense comparison of the two EL records (field type ${\rm E}$ in one
chamber coordinate and ${\rm L}$ in the other; \cref{def:dense-comparison},
and \cref{subsec:push:layer72spectral} for the codes), and the tail law with its one-way repairs
(\cref{lem:tail-prices}). Others accept $\mathcal E_6$ but not all of
$\mathcal E_{12}$: the mirror merge repairs and accepting two-layer
constants (\cref{lem:merge-repair,lem:record:accepting-c2}),
the signed fold records of the middle (\cref{def:record:mixture}), the positive mixture of the six LE records (the field types
${\rm L}$ and ${\rm E}$ exchanged; \cref{prop:le-mixture}), and the upper records $G_{01}L$, $G_{10}L$ with their
reserved one-layer companions, which accept $\mathcal E_6$ and are retained
(\cref{lem:six-query-retention}). So $\mathcal E_{12}$ is only an acceptance set
of parts of the construction: the comparison theorem is applied with
$\mathcal E_6$, $k=6$ (what twelve queries would need is in
\cref{subsec:obs:twelve}). In the density-linear baseline (density exponent $\theta=1$), the exception of the clique prefix is
charged with the two-pair mass $\eta_{\mcl}=2^{1-\mcl}-2^{-2\mcl}$, its exact
mass for $\mathcal E_{12}$ and an upper bound for $\mathcal E_6$
(\cref{lem:frontier:iidlocal}(v)).
\end{remark}

\subsection{Which graph law is used in each prime range}\label{subsec:design-laws}

The construction is one set of comparison data in the sense of Part~III. Its
objects are built in the sections named below; here they are listed together.

\begin{definition}[The construction]\label{def:construction}
The comparison data of Part~V, at the parameters of
\cref{def:record:parameters}, are the following.
\begin{enumerate}[label=(\alph*)]
\item \emph{Frame data} (\cref{def:cp0}). The slot array $V$ with $\mathfrak F$
 and $\iota$. At scale $N$ the atoms of \cref{def:record:functional}: the
 prefix atoms, the blocks of the $33$ middle pieces and of the $310$ upper
 bands, and a tail atom $\{p\}$ for every prime $\Pfour\le p\le N^{\beta}$. The
 queries $\mathcal E_6$. The index law: independent indices
 $\omega_e=(b_e,s_e)\in\Z/\Kint\Z\times\Z/\Mzero\Z$, $e\in\mathcal E_6$, with
 coupled index $(b_e+1,s_e+r_e)$, where each prime-power component of
 $r_e$ is the square of a uniform residue conditioned so that the residue
 itself has $p$-adic valuation less than $64$. The independent depths $J_p$ have tails
 $\Pr(J_p\ge j)=\delta_{p,j}$ of \cref{lem:imp:padic-envelope}; the interval
 coordinate $b_e$ is always revealed (conditioned on, not averaged over;
 \cref{def:hyb:coupled,prop:imp:padic-price}).
\item \emph{Reference} $\mu=\bigotimes_{\Blk}\mu_{\Blk}$ (\cref{def:cp1}): at a
 prefix atom the square-clique law (\cref{def:clique-law}); on a middle block
 the twelve-label five-type completion of \cref{lem:front12:reference}, a law
 on one chamber coordinate and the layer pulled back to $V$, whose two-layer
 root is the uniform labelling law on the pieces below $2^{30}$ and the
 positive root of \cref{def:positive-root} on the seven pieces from $2^{30}$;
 on an upper block the $25$-state five-type completion of the two-layer graph
 law, with independent one-layer fold weights (\cref{def:record:upper-bands},
 \cref{lem:upper-reference-rp}); at a tail atom the tail law of
 \cref{def:tail-law} at the single prime.
\item \emph{First stage and bracket} (\cref{def:cp2,def:cp3}). The first-stage
 atoms are the middle blocks. There the records $F_{\rm fold}$ (the two
 nonaccepting one-layer fold components on the reserved diagonal), $E_{\rm fold}$ (the two
 reserved-folded matching components), $E_{\rm add}$, the polarization
 correction $G$, the reserved half diagonals $H$, and the one-stage repair
 $E_{\rm new}$ of $E_{\rm fold}+E_{\rm add}$ give, with $\rho$ the intermediate law obtained from $\mu$ by polarization before the one-stage repair, and $\Nbrz$ the auxiliary dominator of \cref{def:cp3},
 \begin{gather*}
  \rho=\mu-F_{\rm fold}+E_{\rm add}+G,\qquad
  \lambda=\rho+\delta_E,\qquad \delta_E=E_{\rm new}-E_{\rm fold}-E_{\rm add},\\
  \Nbr=\lambda+F_{\rm fold}+H,\qquad \Nbrz=\mu+E_{\rm add}+G+H
 \end{gather*}
 (\cref{def:record:frame}). The record differences of $\delta_E$ of type $E_0$, whose
 sharp (O2) bound is at their first-coordinate cut $\vartheta_1$, form one of
 the two tags of (CP2)(b), and those of type $E_1$, whose sharp bound is at
 $\vartheta_2$, the other. Off the middle blocks,
 $\lambda_{\Blk}=\Nbr_{\Blk}=\mu_{\Blk}$.
\item \emph{Comparison law} $\nu=\bigotimes_{\Blk}\nu_{\Blk}$ (\cref{def:cp4}),
 with the normalizing masses of $\mu$, as described in the next paragraph. On a
 middle block $\nu_{\Blk}$ is a finite real combination of nonnegative measures
 that accept the queries, each of the two signed fold records contributing
 three product components (\cref{def:record:mixture,lem:middle-comparison-prices});
 every other $\nu_{\Blk}$ is nonnegative. As the middle blocks are finitely
 many, the variation factor of \cref{def:cp4} is a fixed finite number.
\item \emph{Classes, envelopes and the mean credit}
 (\cref{def:cp5,def:cp6,def:cp7}). The variance classes of degrees one to four
 with their gains and positive backgrounds (\cref{sec:variance}); the positive
 joint envelope classes with $n\ge1$ fold factors and $m$ matching
 factors in the tensor product over distinct middle blocks defining the class,
 $n+m\in\{2,4,6\}$, the aligned and crossed matching pairs and the matching
 classes of degrees four and six, with the class bounds of
 \cref{cor:envelope-costs} in the auxiliary prices of
 \cref{lem:middle-comparison-prices}, the invariant background at its assigned
 tag price $\Uinv$ (\cref{lem:oct24:invariant-background}) and the envelope
 portfolio of \cref{tab:data:envelopes}
 (\cref{thm:envelopes,lem:class-credit,lem:orbit-coupon});
 matching envelopes whose sides are the larger of the signed direct-Gram side
 and the auxiliary side, rounded up, with chosen side diagonals
 (\cref{prop:record:matching,prop:signed-sides}); and the mean credit
 (\cref{prop:record:mean-credit}).
\end{enumerate}
\end{definition}

\paragraph{The comparison law in words.}
On a middle block the comparison law keeps the two-layer root and the aligned
fold $F_0$, which accept the queries. It replaces the two other two-layer fold
records $F_1,F_2$ by the signed whole-block law: on the block, the combination
$\bigotimes_pP_{++,p}+(\tlaw/2)\bigotimes_pP_{-+,p}-(\tlaw/2)\bigotimes_pP_{--,p}$
of three positive accepting product laws, each of which is at every prime the
joint law (Markov coupling) of two nonnegative kernels, both
forcing the extra cross edge, whose two marginals are the actual
half marginals of the reference fold record (an \emph{accepting
Markov coupling})
(\cref{def:record:mixture,lem:signed-components}). The combination is taken
after the block products, so its total variation is at most $1+\tlaw$ whatever
the number of primes; $F_2$ is the reflected transpose of $F_1$, and the case
$\tlaw=0$ is the mirror pairing of \cref{lem:mirror}. It replaces the two
two-layer matching records by the exact-marginal law $K_1$
(\cref{lem:record:exact-marginal-e}), repairs every two-layer merge record
by its actual-marginal mirror law (\cref{lem:merge-repair}), replaces
the two-layer constants by the accepting edge law
(\cref{lem:record:accepting-c2}), and, relative to the first-stage law
$\lambda$, deletes the folded merge and constant
records and the polarization correction $G$,
and keeps $E_{\rm new}$ unchanged. On an upper block it keeps every record that accepts
the six queries, among them the records $G_{01}L$, $G_{10}L$ (the two EL records of state
$RL$) and their reserved
one-layer companions (\cref{lem:six-query-retention}); it replaces the six
degree-two LE records by the positive mixture of four separately repaired
six-query laws (\cref{prop:le-mixture}), the two remaining degree-two EL records
by the dense positive comparison
(\cref{def:dense-comparison,prop:dense-comparison}), each of the other $84$
nonaccepting two-layer records of degrees two and three by its one-way repair,
and deletes every other nonaccepting record (\cref{def:upper-comparison}); the
tail atoms use the one-way replacements of the tail law at single primes, with
the all-prime caps (\cref{lem:tail-prices}). At each prefix prime it deletes
the exception of \cref{lem:frontier:iidlocal}. Every change is either a repair
with the same actual half marginals at $\iota$ as the record it replaces, or
the deletion of a positive record. Every component of the comparison law is
nonnegative and accepts the six queries, and only the middle blocks carry
signed combinations. The comparison law is not reflection positive; it is
compared with the first stage only at $\iota$, through the bracket $\Nbr$ and
the norm envelopes of its changes, with the symmetric and skew parts priced
separately.

\paragraph{One recipe per range.}
\Cref{tab:ranges} lists, for each range of primes, the reference law and the
reason for its positivity, the change made by the comparison law, and where
the range enters the closing number $\Theta$.

{\small
\setlength{\LTpre}{6pt}\setlength{\LTpost}{6pt}
\begin{longtable}{@{}>{\raggedright\arraybackslash}p{2.55cm}%
>{\raggedright\arraybackslash}p{4.55cm}%
>{\raggedright\arraybackslash}p{3.3cm}%
>{\raggedright\arraybackslash}p{4.25cm}@{}}
\caption{The ranges of primes.}\label{tab:ranges}\\
\toprule
Range & Reference law and its positivity & Comparison change & Contribution to $\Theta$\\
\midrule
\endfirsthead
\toprule
Range & Reference law and its positivity & Comparison change & Contribution to $\Theta$\\
\midrule
\endhead
\bottomrule
\endfoot
Small modulus: $2^{129}$ and $p^{128}$ for $3\le p<Q_3$, $p\equiv3\ (4)$ (not atoms)
& no law; an external index per query, with coupled index $(b+1,s+r)$;
 doubly stochastic kernels, no positivity needed
& none
& the conductor moment $R_I=M_{\rm pad}(\Kint)$ in $J$ and $W$, at $\Kint=\PushK$
 (\cref{prop:oct24-index-cost,lem:index-moments})\\
\addlinespace
Prefix: $p\equiv1\ (4)$, $p<\Pfour$; atoms $\{p\}$, $e_p=2\mcl$
& independent draws from a random square clique; at the chamber
 cuts of the first coordinate the square of a conditional mean,
 the other cuts fix it samplewise
 (\cref{lem:frontier:iidlocal})
& delete the zero-difference exception
& mass at most $\eta_{\mcl}$; total price below $2^{513-\mcl+1024\mcl a}$, in
 $S$ and $T$ only, inside the allowance $10^{-12}$
 (\cref{lem:frontier:iidenergy})\\
\addlinespace
Middle: $p\equiv3\ (4)$ in $[Q_3,\Pthree)$; $33$ pieces, blocks with
 $\log q_{\Blk}\le L$
& twelve-label five-type completion on one chamber coordinate and the layer
 (\cref{lem:front12:reference}), its root uniform, or positive on the seven
 pieces from $2^{30}$ (\cref{def:positive-root}); chamber cuts by the
 completion lemma with the caps of the root
 (\cref{lem:layer12-caps,prop:l20-collisions,prop:root-chamber,prop:root-layer})
 and the rows, the layer flip by fold pairing (\cref{lem:five-type})
& first stage: polarized folds, one-stage repair $E_{\rm new}$ in two tags;
 comparison: signed whole-block $F_1,F_2$, exact-marginal $E$, mirror
 merges, accepting two-layer constants, folded deletions
& $L_0$ and the energy factor $B$ through the tagged series
(\cref{sec:middle-frame}); the direct
 prices of the middle records in $S$, $A$, and the auxiliary prices of the
 fold and matching records in the class bounds
 (\cref{lem:middle-comparison-prices}); variance and envelope credits in
 $\mathcal E$ (\cref{sec:variance}, \cref{cor:envelope-costs}); the prices of
 the matching and other records and the side excess $T_K$
(\cref{prop:signed-sides}(iii)) of the matching
 envelopes in the total $T$, the mean savings in $m$ (\cref{sec:matching-mean})\\
\addlinespace
Upper: $p\equiv3\ (4)$ in $[\Pthree,\Pfour)$; $310$ bands
& $25$-state two-layer five-type completion with independent reserved fold
 weights and observed reserved caps
 (\cref{lem:upper-reference-rp,prop:observed-caps}); root caps by the direct
 Route~B criterion below $2^{52}$, with graded counts and interpolated
 coefficients, and by the Route~A criterion from $2^{52}$
 (\cref{thm:record:direct-star,thm:spectral-criterion}); joint fold and
 matching rows (\cref{lem:five-type})
& one-way repairs of $84$ records, a dense comparison of two EL records, a
 positive mixture of the six LE records; deletion of the other nonaccepting
 records
& $S_{\rm u}(a)$, $A_{\rm u}(a)$ in $S$, $A$ and $T$, the band ledger repriced
 band by band (\cref{prop:record:upper})\\
\addlinespace
Tail: every $p\ge\Pfour$; atoms $\{p\}$, $e_p=1$
& the tail law at one prime (\cref{def:tail-law}), at class-one primes with
 the all-prime caps; mixing \ref{it:record:R4} for $p\ge\Pfive$
 (\cref{lem:tail-law,lem:record:tail-mixing})
& one-way repairs of the $94$ records that do not accept the twelve queries,
 with the all-prime caps; deletion of the other nonaccepting records
& per prime at most $Cp^{-3/2+2a}$ with the absolute constant $C$
  of \cref{lem:tail-prices}, inside $S_{\rm u}(a)$ and
  $A_{\rm u}(a)$
 (\cref{lem:tail-prices})\\
\end{longtable}
}

\subsection{Three uses of ``admissible'', and why each holds}\label{subsec:design-admissibility}

The energy estimates of Part~II apply to the energies of the comparison when
the laws and the reductions are admissible (\cref{def:admissible}), and (CP8)(a)
asks this of every law of the construction. Each construction lemma of
\cref{sec:frontier:iidprefix} and of the sections after it states the
admissibility of its laws; \cref{tab:admissibility} collects them, with the two
facts that give (O1) and (O2) for each.

{\small
\setlength{\LTpre}{6pt}\setlength{\LTpost}{6pt}
\begin{longtable}{@{}>{\raggedright\arraybackslash}p{3.45cm}%
>{\raggedright\arraybackslash}p{4.4cm}%
>{\raggedright\arraybackslash}p{3.65cm}%
>{\raggedright\arraybackslash}p{3.15cm}@{}}
\caption{The laws of the construction and their admissibility.}\label{tab:admissibility}\\
\toprule
Law & Sign; one-slot marginals (O1) & Sharp bounds (O2) & Proved in\\
\midrule
\endfirsthead
\toprule
Law & Sign; one-slot marginals (O1) & Sharp bounds (O2) & Proved in\\
\midrule
\endhead
\bottomrule
\endfoot
prefix: $\lambda_p$, the exception $\lambda^{\rm ex}_p$, $\lambda_p-\lambda^{\rm ex}_p$
& nonnegative; invariant under $\mathfrak G_p$
& none; the exception is charged by its mass
& \cref{lem:frontier:iidlocal}(vi)\\
middle and upper references $\mu_{\Blk}$
& nonnegative; invariant under the maps of \ref{it:record:R3}
& none
& \cref{lem:front12:reference}, \cref{lem:positive-root-properties,lem:imp:coherent39-law},
 \cref{lem:upper-reference-rp}\\
first-stage and bracket records $F_{\rm fold}$, $E_{\rm fold}$, $E_{\rm add}$,
 $G$, $H$; $\Nbrz-\mu=E_{\rm add}+G+H$
& nonnegative, reserved diagonal; images of invariant laws under vertex maps,
 then diagonal copies
& none; variation and mass prices
& \cref{def:record:frame}, \cref{lem:record:gram}\\
one-stage repair $E_{\rm new}$
& nonnegative, reserved diagonal; invariant
& $k_{\rm new}(X)b^{-1/2}$, and $s_{\rm fold}(X)/b$ for the symmetric part, for
 every value of the shared label, where $X$ is the piece start and $b$ the
least prime of the block
& \cref{lem:record:one-stage-folded-e}\\
middle comparison records: the three product components of each signed fold
 record, the law $K_1$, merge and two-layer constant repairs
& nonnegative; invariant under the maps of \ref{it:record:R3}, independently
 at the primes of the block; the actual half marginals of the replaced record
& operator norms between actual half spaces at $\iota$, nothing conditioned
& \cref{lem:signed-components}, \cref{prop:signed-bounds},
 \cref{lem:record:exact-marginal-e}, \cref{lem:merge-repair,lem:record:accepting-c2};
 \cref{lem:middle-comparison-prices}\\
upper and tail comparison laws: one-way repairs, the dense EL comparison, the
 components of the LE mixture
& nonnegative; the actual half marginals of the replaced record, each mixture
 component repaired to them before mixing
& as for the middle records
& \cref{prop:record:upper}, \cref{prop:dense-comparison},
 \cref{prop:le-mixture}, \cref{lem:tail-prices}\\
deleted records (middle, upper, tail)
& records of the reference
& none; charged by mass (\cref{rem:cp:charging})
& \cref{lem:middle-comparison-prices}, \cref{prop:record:upper},
 \cref{lem:tail-prices}\\
tail law $\mu_p$, $p\ge\Pfour$
& as the upper references, at one prime
& none
& \cref{lem:tail-law}\\
selected laws of the variance credits
& positive; invariant under independent translations at the primes of their
 blocks
& $d=1$ for a probability coupling; centred norms of a record's own operator at
 $\iota$
& \cref{def:record:eligible}, \cref{lem:var-gain}\\
couplings of the positive joint envelopes
& normalized positive couplings, among them the singleton couplings
 $M_{ij}=V_iV_j^*$ (a uniform label and independent conditioned halves);
 invariant under translations at each prime
& none (the envelope bounds are brackets at $\iota$ between actual half spaces)
& \cref{thm:envelopes}(iv), \cref{lem:oct24:mask-schur}\\
matching-envelope diagonals $M_{s,\Blk}-\Nbr_{\Blk}$
& positive, reserved diagonal (norm envelopes and chosen side diagonals);
 normalizations with uniform one-slot marginals
& block Grams at $\iota$ between actual half spaces
& \cref{prop:record:matching}, \cref{prop:signed-sides}\\
companions of the mean credit
& positive normalized laws fixed by $\iota$, uniform one-slot marginals
& none
& \cref{lem:record:companion-energy}\\
\end{longtable}
}

\paragraph{How the table is read.}
For (O1), every law of the table is nonnegative (a signed fold record of the
middle enters only through its three nonnegative product components, each of
which is such a law), and its one-slot marginals are uniform for one of two
reasons named in the table: it is invariant under the translations, or all
square-affine maps, at the primes of its atom, so that
\cref{lem:marginals}\ref{it:marg:one} applies; or it is nonnegative with the
actual half marginals of such an invariant law, whose one-slot marginals it
therefore shares. Invariance passes to every law that the reductions of
Part~III produce from these: to variations, normalizations, restrictions to
invariant events, half marginals and their diagonal copies, images under folds
and pullbacks, products over disjoint atoms and mixtures
(\cref{lem:marginals}\ref{it:marg:closure}). Every law of the table is a law of
the field labels of its atoms, and the index is external (\cref{def:cp0}(d)); so
no law depends on the interval index $b_e$, the residue $s_e$, or the
inclusion indices of \cref{def:hyb:coupled}, which is the last
clause of (O1). For (O2), the third column says why every sharp bound applied
to the law holds uniformly in the shared labels: either it is an operator norm
between actual half spaces with nothing conditioned, or it is proved for every
value of the shared label. The remaining two obligations concern the
reductions, not the laws. (O3): the index-support projections are taken before
the field folds and the inclusion conductors are averaged last, in the
reductions of Part~III (\cref{lem:record:common-mean,lem:record:wrap-expansion,lem:envelope-energy}),
and no section of Part~V changes this order. (O4): every chamber fold commutes
with $\iota$ on the slot array (\cref{lem:slot-array}(iii)).

\subsection{Where the hypotheses are verified (guide to Part~V)}\label{subsec:design-table}

\Cref{tab:hypotheses} lists each hypothesis of \cref{thm:comparison}, what
Part~V shows for the construction of \cref{def:construction}, and the statements
that show it. The frame data (CP0) are verified by the parameters above, by the
slot array (\cref{lem:slot-array}) and by the index law
(\cref{prop:oct24-index-cost,lem:index-moments}); the hypotheses (CP1)--(CP7)
range by range in the following sections; the admissibility clause (CP8)(a) by
\cref{tab:admissibility}; and the summability clause (CP8)(c) because only the
tail atoms are infinitely many and their prices are summable
(\cref{lem:tail-prices}(iv)), while the auxiliary prices of (CP5) are nonzero
only at the finitely many middle atoms. The robustness clause (CP8)(b) follows
from \cref{lem:record:rec}(i) (\cref{rem:cp:robust-final}):
under the inflation of \cref{def:robust} (multiplying structured
energy bounds by $1+z$ and adding remainders at most $zY$), every
certified inequality (gate) whose left side changes is strict,
except two: the orbit saving of
\cref{lem:orbit-coupon}, which is a nonnegativity condition on a
coefficient that only grows under inflation, and the hypothesis
$\chi<1$ selecting the branch of the restoration function
\eqref{eq:density:restoration}, which inflation preserves. The
remaining non-strict gates concern fixed numbers of the
construction. Only the
finitely many middle blocks carry signed comparison laws
(\cref{def:construction}(d)), so the variation factor of \cref{def:cp4} is
finite and independent of the scale.

\paragraph{Consequences and what to read next.}
This section uses the slot array and its queries (\cref{def:slot-array},
\cref{lem:slot-array}(iii),
(iv), (v)); the comparison data and the hypotheses (CP0)--(CP8)
(\cref{def:cp0,def:cp1,def:cp2,def:cp3,def:cp4,def:cp5,def:cp6,def:cp7,def:cp8});
admissibility and its closure under the operations of the reductions
(\cref{def:admissible}, \cref{lem:marginals}); the functional, the analytic stage
and the repricing (\cref{def:record:functional}, \cref{def:record:analytic-stage},
\cref{cert:record:analytic-stage}, \cref{lem:record:repricing},
\cref{cor:first-stage-repricing}); and, as pointers, the construction lemmas named
in \cref{tab:ranges,tab:admissibility}.

\paragraph{Hypotheses verified in this section (\cref{tab:hypotheses}).}
(CP0)(b) for the construction: the atoms and exponents of
\cref{def:record:parameters}, with $a<\sigma$, the stage inequalities (S1)--(S7)
and the analytic gates on $0<a\le9/10000$ (\cref{cert:record:analytic-stage});
(CP0)(c): the six queries have disjoint endpoints, with sources and reflected
targets in $V_+$ (\cref{subsec:front6:construction}, \cref{lem:slot-array}(iv));
for (CP4), that only the finitely many middle blocks carry signed combinations,
so that the variation factor is finite and independent of the scale
(\cref{def:construction}(d)); and the summary of (CP8)(a) in
\cref{tab:admissibility}. These are the rows ``(CP0)(b)'', ``(CP0)(c)'' and
``(CP8)(a)'' of \cref{tab:hypotheses}, and the finiteness of the variation
factor in the row ``(CP4) middle $F$'', whose variation $1+\tlaw$ per block is
that of the signed fold record (\cref{def:record:mixture}(c)).
